\begin{ouverture}[grandeurs]{14}{Horaires et durées}{Partie 3 · Grandeurs et mesures}
{Heures · Minutes · Instants · Temps écoulé}
\objectif{Lire l’heure sur une horloge à aiguilles.}
\objectif{Convertir des durées entre heures et minutes.}
\objectif{Calculer un début, une fin ou une durée.}
\prerequis{Division euclidienne et addition posée.}
\end{ouverture}
\pagestyle{otMain}\essentiel\fontsize{11}{15}\selectfont
\phantomsection\label{cours-14-1}
\essSection{1}{Lire l’heure}
\begin{definition}Sur une horloge à aiguilles, la petite aiguille indique les heures, la grande les minutes. Chaque graduation d’heure vaut 5 minutes pour la grande aiguille : quand elle pointe le 7, il est 35 minutes.\end{definition}
\label{fig-cours-14-1-cours-horloge}\figureGeo{cours-horloge}

\begin{propriete}Un jour dure 24 heures et une heure dure 60 minutes. L’après-midi, on ajoute 12 aux heures du cadran : 2 h de l’après-midi s’écrit 14 h ; 20 h 30 se dit aussi « 8 h 30 du soir ».\end{propriete}
\begin{savoirfaire}[Minutes et heures]Convertir en passant par 60 : $1$ h $30$ min $=60+30=90$ min ; et 180 min $=3\times60$ min $=3$ h. Pour 95 min : $95=60+35$, soit 1 h 35 min.\end{savoirfaire}
\clearpage
\phantomsection\label{cours-14-2}
\essSection{2}{Calculer des durées}
\begin{savoirfaire}[Avancer par étapes]Pour la durée entre 9 h 40 et 10 h 15, avancer par étapes : de 9 h 40 à 10 h, 20 min ; de 10 h à 10 h 15, 15 min. Total : 35 min. Une frise des instants aide à poser les étapes.\end{savoirfaire}
\begin{exemple}Heure de fin : partir de 10 h 20 et durer 45 min. $10$ h $20+40$ min $=11$ h, puis encore 5 min : 11 h 05. Heure de début : finir à 12 h 05 après 30 min, c’est partir à 11 h 35.\end{exemple}
\begin{attention}Un horaire est un instant ; une durée est un temps écoulé. On n’additionne pas deux horaires. Dans les calculs, 60 minutes s’échangent contre une heure : il n’existe pas d’horaire à 70 minutes.\end{attention}
\begin{propriete}Pour comparer des durées, les exprimer dans la même unité : 1 h 15 min $=75$ min est plus court que 80 min.\end{propriete}

\clearpage\rubrique{Découverte}
\begin{blocLivre}{ex-14-1}\begin{exercice}[1]{Lire deux horloges}
Lire l’heure affichée par chacune des deux horloges, le matin.
\label{fig-ex-14-1}\figureGeo{hor-1}
\end{exercice}\end{blocLivre}
\begin{blocLivre}{ex-14-2}\begin{exercice}[1]{Placer les aiguilles}
Dessiner les aiguilles d’une horloge indiquant 16 h 45.
\label{fig-ex-14-2}\figureGeo{hor-2}
\end{exercice}\end{blocLivre}
\begin{blocLivre}{ex-14-3}\begin{exercice}[1]{Matin ou après-midi}
Écrire en heures de l’après-midi ou du soir : 15 h ; 20 h 30. Écrire avec les heures du jour : « 5 h du soir ».
\end{exercice}\end{blocLivre}
\begin{blocLivre}{ex-14-4}\begin{exercice}[1]{Convertir}
Convertir : 2 h en minutes ; 1 h 30 min en minutes ; 180 min en heures.
\end{exercice}\end{blocLivre}
\begin{blocLivre}{ex-14-5}\begin{exercice}[1]{Dans la même heure}
Un cours commence à 14 h 10 et finit à 14 h 45. Combien de temps dure-t-il ?
\end{exercice}\end{blocLivre}

\rubrique{Entraînement}
\begin{blocLivre}{ex-14-6}\begin{exercice}[2]{Franchir l’heure}
Quelle durée s’écoule entre 9 h 40 et 10 h 15 ?
\end{exercice}\end{blocLivre}
\begin{blocLivre}{ex-14-7}\begin{exercice}[2]{L’heure de fin}
Un entraînement commence à 10 h 20 et dure 45 min. À quelle heure finit-il ?
\end{exercice}\end{blocLivre}
\begin{blocLivre}{ex-14-8}\begin{exercice}[2]{L’heure de départ}
Un trajet de 30 min se termine à 12 h 05. À quelle heure a-t-il commencé ?
\end{exercice}\end{blocLivre}
\begin{blocLivre}{ex-14-9}\begin{exercice}[2]{Des minutes aux heures}
Convertir 95 min en heures et minutes, puis 2 h 05 min en minutes.
\end{exercice}\end{blocLivre}
\begin{blocLivre}{ex-14-10}\begin{exercice}[2]{La plus longue durée}
Qui écoute le plus de musique : Ana pendant 1 h 15 min ou Bob pendant 80 min ?
\end{exercice}\end{blocLivre}

\rubrique{Problèmes}
\begin{blocLivre}{ex-14-11}\begin{exercice}[3]{La séance de cinéma}
Un film commence à 14 h 50 et dure 1 h 45 min. À quelle heure se termine-t-il ?
\end{exercice}\end{blocLivre}
\begin{blocLivre}{ex-14-12}\begin{exercice}[3]{La journée de classe}
La classe travaille de 8 h 30 à 11 h 45, puis de 13 h 30 à 16 h 30. Calculer la durée totale de classe.
\label{fig-ex-14-12}\figureGeo{hor-12}
\end{exercice}\end{blocLivre}
\begin{blocLivre}{ex-14-13}\begin{exercice}[3]{Le train de 9 h 48}
Un train part à 9 h 48 et arrive à 11 h 23. Combien de temps dure le voyage ?
\end{exercice}\end{blocLivre}
\begin{blocLivre}{ex-14-14}\begin{exercice}[3]{Les départs du bus}
Un bus part toutes les 20 minutes ; le premier départ est à 7 h. Écrire les horaires de départ jusqu’à 8 h inclus. Combien y a-t-il de départs entre 7 h et 9 h, les deux horaires compris ?
\end{exercice}\end{blocLivre}
\begin{blocLivre}{ex-14-15}\begin{exercice}[3]{Le gâteau au four}
Un gâteau doit cuire 55 min. Il est enfourné à 17 h 40 et sorti à 18 h 30. Est-il resté assez longtemps au four ?
\end{exercice}\end{blocLivre}

\clearpage
